% =====================================================================
%  B201 Computer Science Lab — Individual Project Report
%  Uygar Tuna Oflas
%  Compile with: pdflatex report.tex  (run twice for the ToC)
% =====================================================================
\documentclass[11pt,a4paper]{article}

\usepackage[T1]{fontenc}
\usepackage[utf8]{inputenc}
\usepackage[margin=2.4cm]{geometry}
\usepackage{xcolor}
\usepackage{titlesec}
\usepackage{enumitem}
\usepackage{hyperref}
\usepackage[round,authoryear]{natbib}
\usepackage{listings}
\usepackage{parskip}

\definecolor{accent}{HTML}{2C44D6}
\definecolor{ink}{HTML}{15181E}
\definecolor{codebg}{HTML}{F4F5F8}

\hypersetup{colorlinks=true, urlcolor=accent, linkcolor=accent, citecolor=accent}

\titleformat{\section}{\large\bfseries\color{ink}}{\thesection}{0.6em}{}
\titleformat{\subsection}{\normalsize\bfseries\color{accent}}{\thesubsection}{0.6em}{}
\titlespacing{\section}{0pt}{14pt}{4pt}

\lstset{
  basicstyle=\ttfamily\footnotesize,
  backgroundcolor=\color{codebg},
  frame=single, rulecolor=\color{codebg},
  breaklines=true, xleftmargin=6pt, framexleftmargin=6pt,
  showstringspaces=false,
}

\title{\vspace{-1.2cm}\textbf{Designing and Building a Computer Science Portfolio}}
\author{Uygar Tuna Oflas \\ \small B201 Computer Science Lab — Gisma University of Applied Sciences}
\date{\small\today}

\begin{document}
\maketitle
\vspace{-0.6cm}

% ---------------------------------------------------------------------
\begin{center}\small
\textbf{Live site:} \url{https://uygartunaoflas.github.io/cs-portfolio/} \quad\textbar\quad
\textbf{Repository:} \url{https://github.com/uygartunaoflas/cs-portfolio}
\end{center}
\vspace{0.2cm}

% ---------------------------------------------------------------------
\section{Introduction}
This report documents the design, implementation and content of a personal
computer-science portfolio, produced for the B201 Computer Science Lab. The brief
frames the task realistically: a hiring manager for a working-student position has
asked for a portfolio that demonstrates technical skill. The deliverable is therefore
two connected things --- a portfolio website hosted on GitHub Pages, and the
repository behind it --- both of which should be clear, well structured and honest
about what they show.

My guiding principle was \emph{restraint with intent}. As an early-career student, I
do not have a long list of large projects, so I chose to present a small number of
focused exercises, each demonstrating one idea cleanly, rather than padding the site.
The report explains how the site and repository are organised, justifies the main
design decisions, lists the tools used, and reflects critically on the result.

% ---------------------------------------------------------------------
\section{Structure of the Portfolio}

\subsection{Website}
The website is a single page divided into labelled sections that a recruiter can scan
in order: a \textbf{hero} introduction, \textbf{skills}, an \textbf{interactive demo},
\textbf{projects}, \textbf{experience}, \textbf{education} and \textbf{contact}. A
single-page layout suits a portfolio of this size: the visitor sees everything by
scrolling, and a sticky navigation bar provides quick jumps without page reloads. This
follows the common guidance that navigation should let users locate information with
minimal effort \citep{nielsen2020}.

The hero states who I am and what I am looking for in one sentence, because the top of
a page carries the most weight in how quickly visitors form a judgement
\citep{lidwell2010}. Below it, four short ``fact'' chips (location, studies, languages,
interests) give context at a glance.

\subsection{Repository}
The repository mirrors the site's content but is organised for reading code rather than
browsing:

\begin{lstlisting}
cs-portfolio/
  index.html            # the website (HTML + CSS + JS in one file)
  cv/                   # LaTeX CV source and compiled PDF
  projects/             # one folder per technical exercise
    knight_path/        # Python BFS + unit tests
    text_insights/      # Python CLI
    guestbook/          # PHP endpoint
  report/               # this report (LaTeX + PDF)
  README.md             # overview and run instructions
\end{lstlisting}

Each project lives in its own folder so it can be read in isolation, and the
\texttt{README.md} acts as a map with run instructions. Keeping a clear, predictable
project layout is a widely recommended habit because it lowers the effort needed to
understand a codebase \citep{wilson2017}.

% ---------------------------------------------------------------------
\section{Design Decisions}

\subsection{A single static page, no framework}
I deliberately built the site with plain HTML, CSS and JavaScript rather than a
framework such as React. For a static portfolio with no dynamic data, a framework would
add a build step and dependencies without adding value; the simplest tool that does the
job is the right one. GitHub Pages serves static files directly \citep{githubpages2024},
so a single \texttt{index.html} deploys with no configuration. Inlining the CSS and
JavaScript keeps the whole site in one reviewable file, which is reasonable at this
scale, though I note the trade-off in the reflection.

\subsection{Visual identity}
Most developer portfolios look alike, so I made specific choices rather than accepting
defaults. The palette is a cool light-grey paper (\texttt{\#ECEFF3}) with a single
indigo accent (\texttt{\#2C44D6}); limiting the design to one accent colour keeps
attention on the content and avoids visual noise \citep{lidwell2010}. The type pairing
is intentional: \emph{Space Grotesk} for headings gives a technical character,
\emph{Inter} for body text is highly legible at small sizes, and \emph{JetBrains Mono}
is used only for labels and code, where a monospace face signals ``this is technical''.
Pairing a display face with a separate, readable body face is a standard typographic
principle \citep{butterick2019}.

\subsection{The interactive demo as a signature}
Rather than a static screenshot, the centre of the page is a working
\textbf{knight's-shortest-path} board. Clicking a start square and a target square runs
a breadth-first search in the browser and draws the optimal route. I chose this for
three reasons: it ties to a personal interest (chess), it shows a real algorithm rather
than describing one, and it connects directly to the Python implementation in the
repository, so the demo and the code reinforce each other. A small interactive moment
also demonstrates DOM manipulation and event handling without a framework.

\subsection{Accessibility and responsiveness}
The site is built to a basic quality floor. Layouts use responsive CSS grids that
collapse to a single column on narrow screens, so the page is usable on a phone. Colour
contrast between text and background was kept high for readability, in line with WCAG
guidance \citep{wcag2018}. Interactive controls are real \texttt{<button>} elements with
\texttt{aria-label}s and visible keyboard focus styles, and all motion is disabled when
the visitor's system requests reduced motion (\texttt{prefers-reduced-motion}). These
are small additions but they reflect the principle that an interface should work for as
many people as possible \citep{wcag2018}.

\subsection{The CV in LaTeX}
The brief requires a LaTeX CV, and LaTeX is a good fit: it produces consistent,
typeset PDFs and separates content from formatting \citep{lamport1994}. I wrote a
custom one-column layout using only standard packages, so it compiles anywhere without
extra installation, and reused the same indigo accent as the website so the CV and site
read as one identity.

% ---------------------------------------------------------------------
\section{Tools and Technologies}
\begin{itemize}[leftmargin=1.3em, itemsep=2pt]
  \item \textbf{HTML, CSS, JavaScript} --- the website, with no framework or build step.
  \item \textbf{GitHub Pages} --- free static hosting served from the repository \citep{githubpages2024}.
  \item \textbf{Git \& GitHub} --- version control and the public repository.
  \item \textbf{LaTeX} --- the CV and this report.
  \item \textbf{Python} --- the knight-path finder (with \texttt{unittest}) and the text-insights CLI.
  \item \textbf{PHP} --- the guestbook endpoint sample.
  \item \textbf{Google Fonts} --- Space Grotesk, Inter and JetBrains Mono.
\end{itemize}

% ---------------------------------------------------------------------
\section{Implementation of the Technical Exercises}

\subsection{Knight Path Finder (Python + JavaScript)}
The core exercise finds the fewest knight moves between two squares. It models the board
as an implicit graph where each square connects to the squares a knight can jump to, and
runs a breadth-first search. BFS is the right choice because it explores the graph in
order of distance, so the first time it reaches the target it has used the minimum number
of moves \citep{cormen2009}. The Python version is the reference implementation and is
covered by unit tests, including the known result that a corner-to-corner trip (a1 to h8)
takes six moves. The JavaScript version on the website uses the same algorithm so the
visible demo and the tested code agree.

\subsection{Text Insights CLI (Python)}
A command-line tool that reads a text file and reports word and sentence counts, average
lengths, and the most frequent content words after removing common stop words. It is a
small exercise in parsing, using dictionaries (via \texttt{collections.Counter}) and
presenting tidy output --- everyday tasks that appear constantly in real work.

\subsection{Guestbook API (PHP)}
A minimal endpoint that accepts short messages and stores them to a JSON file. The point
of the exercise is not the storage but the discipline around input: every field is
checked for presence, type and length, and text is escaped before it is stored, which
guards against the kind of cross-site-scripting issue that arises when user input is
trusted. Validating and encoding untrusted input is a basic security expectation
\citep{owasp2021}. PHP runs on a server rather than on GitHub Pages, so this file is
included as a readable sample rather than a live feature.

% ---------------------------------------------------------------------
\section{Reflection: Strengths, Weaknesses and Future Work}

\subsection{Strengths}
The portfolio is honest and coherent. It does not overstate experience; instead it shows
a small set of working, documented exercises, and the interactive demo proves a claim
rather than asserting it. The site and CV share one visual identity, the repository is
tidy and explained, and the whole thing deploys with no build step, which keeps it easy
to maintain.

\subsection{Weaknesses}
There are real limits. Inlining the CSS and JavaScript into one HTML file is convenient
at this size but would become hard to maintain as the site grows; splitting them into
separate files would be the next step. The project set is deliberately small and leans
toward classic exercises rather than a larger, original application. There are no
automated tests for the front-end JavaScript, and accessibility has been addressed by
hand rather than verified with an audit tool.

\subsection{Future work}
Given more time I would: (1) separate the stylesheet and script and add a light build or
linting step; (2) add a larger original project --- for example a small full-stack
application --- to show work beyond isolated exercises; (3) run an accessibility audit
(e.g.\ Lighthouse) and act on its findings; and (4) add a short ``writing'' or
``notes'' section, since the ability to explain technical work clearly is itself a skill
worth showing. These changes would deepen the portfolio without changing its core
principle: show real, understood work, clearly.

% ---------------------------------------------------------------------
\section*{References}
\begin{thebibliography}{99}
\setlength{\itemsep}{2pt}

\bibitem[Butterick(2019)]{butterick2019}
Butterick, M. (2019) \emph{Practical Typography}. 2nd edn. Available at: \url{https://practicaltypography.com} (Accessed: 26 June 2026).

\bibitem[Cormen et al.(2009)]{cormen2009}
Cormen, T.H., Leiserson, C.E., Rivest, R.L. and Stein, C. (2009) \emph{Introduction to Algorithms}. 3rd edn. Cambridge, MA: MIT Press.

\bibitem[GitHub(2024)]{githubpages2024}
GitHub (2024) \emph{GitHub Pages Documentation}. Available at: \url{https://docs.github.com/en/pages} (Accessed: 26 June 2026).

\bibitem[Lamport(1994)]{lamport1994}
Lamport, L. (1994) \emph{LaTeX: A Document Preparation System}. 2nd edn. Reading, MA: Addison-Wesley.

\bibitem[Lidwell et al.(2010)]{lidwell2010}
Lidwell, W., Holden, K. and Butler, J. (2010) \emph{Universal Principles of Design}. Beverly, MA: Rockport.

\bibitem[Nielsen(2020)]{nielsen2020}
Nielsen, J. (2020) \emph{10 Usability Heuristics for User Interface Design}. Nielsen Norman Group. Available at: \url{https://www.nngroup.com/articles/ten-usability-heuristics/} (Accessed: 26 June 2026).

\bibitem[OWASP(2021)]{owasp2021}
OWASP (2021) \emph{OWASP Top 10: Web Application Security Risks}. Available at: \url{https://owasp.org/www-project-top-ten/} (Accessed: 26 June 2026).

\bibitem[W3C(2018)]{wcag2018}
W3C (2018) \emph{Web Content Accessibility Guidelines (WCAG) 2.1}. Available at: \url{https://www.w3.org/TR/WCAG21/} (Accessed: 26 June 2026).

\bibitem[Wilson et al.(2017)]{wilson2017}
Wilson, G. et al. (2017) `Good enough practices in scientific computing', \emph{PLOS Computational Biology}, 13(6).

\end{thebibliography}

\end{document}
